\documentclass{situationist-preprint}

\title{Galileo's Falling Bodies}
\msid{SCI-15-b}
\msfamily{SCI-15}
\msdiscipline{Philosophy of Science}
\mstype{Companion manuscript}
\mskeywords{thought experiments; Galileo; falling bodies; Aristotle; argument view; Gendler}

\begin{document}
\maketitle

\begin{abstract}
Galileo's thought experiment of the tied bodies is the standard case for the view that some thought experiments give insight beyond argument. SCI-15-a holds that the evidential force of a thought experiment is that of the argument it encodes. We reconstruct Galileo's argument with its premises explicit and show that the reconstruction predicts both what the thought experiment establishes and where it fails: it refutes the thesis that natural speed of fall is fixed by weight alone, and it does not by itself establish that all bodies fall at the same rate, which requires a further premise about resistance. We answer the objection that the thought experiment works by revealing a conceptual tension that no argument captures.
\end{abstract}

\section{The case}

In the \emph{Discourses}, Salviati asks Simplicio to consider a heavy stone and a light stone, tied together, falling \citep{Galileo1638}. On the Aristotelian view that heavier bodies fall faster, the light stone, falling more slowly, should retard the heavy one, so that the pair falls more slowly than the heavy stone alone; but the pair is heavier than the heavy stone, and so should fall faster. The contradiction is taken to show that heavy and light bodies fall at the same rate. \citet{Brown1991} treats the case as a platonic insight into a law of nature; \citet{Norton1996} as an argument in picturesque form.

\section{The encoded argument}

By Definition~2.2 of SCI-15-a, the argument a thought experiment encodes has as premises the propositions on which the thinker's judgment about the scenario depends. For the tied bodies they are:
\begin{enumerate}
\item[(A)] Each body has a natural speed of fall, which is an increasing function of its weight alone.
\item[(M)] When a faster and a slower body are joined, the slower retards the faster and the faster hastens the slower, so the composite has a speed strictly between theirs.
\item[(C)] Two bodies tied together form a body whose weight is the sum of their weights.
\end{enumerate}
From (A) and (C), the composite falls faster than the heavy stone; from (A) and (M), more slowly. Hence (A), (M) and (C) are jointly inconsistent. Since (M) and (C) are accepted by the Aristotelian as well as by Galileo, the argument refutes (A).

\begin{proposition}\label{p:what}
The tied-bodies thought experiment establishes that natural speed of fall is not an increasing function of weight alone. It does not establish that all bodies fall at the same rate.
\end{proposition}

\begin{proof}
The first claim is the conclusion of the reconstructed argument. For the second, the premises are consistent with speeds that depend on properties other than weight, such as shape, size or density relative to the medium, and in air bodies do fall at different rates. That all bodies fall at the same rate follows only with a further premise, that speed of fall depends on nothing but weight in the absence of resistance, which together with the refutation of (A) yields a speed independent of weight. Galileo supplies this premise in the subsequent discussion of fall in media of decreasing density, not in the thought experiment.
\end{proof}

Proposition~\ref{p:what} is what Theorem~3.2 of SCI-15-a predicts: the force of the thought experiment is the force of its argument, no more. Corollary~3.3 of SCI-15-a, that a thought experiment fails exactly when its argument has an unjustified premise or a bad inference, explains why it fails as a demonstration of universal equal fall: the needed premise about resistance is not among its premises. A platonist account, on which the thinker grasps the law directly, does not explain why the insight extends exactly as far as the argument and stops where the argument stops.

\begin{objection}
\citet{Gendler1998} argues that the thought experiment works by bringing the thinker to see that weight, as Aristotle understood it, cannot be the relevant quantity, since a tied pair has no determinate status as one body or two. This reconceptualization is not an inference from premises.
\end{objection}

\begin{reply}
The observation that the tied pair can be regarded as one body or as two is premise (C) together with the recognition that (M) treats it as two. That the speed of the pair should not depend on whether a string joins its parts is a further premise, and a plausible one. Once it is stated, the reasoning is an argument, and its force is the force of that premise. The reconceptualization is the thinker's coming to accept a premise, which Principle~P2 of SCI-15-a allows to be produced by imagining the scene, provided its justification derives from information already possessed, as Principle~P1 requires.
\end{reply}

\section{Conclusion}

The tied bodies refute the Aristotelian law because the argument they encode is valid and its premises are shared; they fail to establish universal equal fall because the argument lacks the needed premise. Both verdicts are those of the argument view.

\begin{thebibliography}{99}\small
\bibitem[Brown(1991)]{Brown1991} Brown, J.~R. (1991). \emph{The Laboratory of the Mind: Thought Experiments in the Natural Sciences}. Routledge.
\bibitem[Galileo(1638)]{Galileo1638} Galilei, G. (1638). \emph{Discorsi e dimostrazioni matematiche intorno a due nuove scienze}. Lodewijk Elzevir, Leiden. First Day.
\bibitem[Gendler(1998)]{Gendler1998} Gendler, T.~S. (1998). Galileo and the indispensability of scientific thought experiment. \emph{The British Journal for the Philosophy of Science}, 49(3), 397--424.
\bibitem[Norton(1996)]{Norton1996} Norton, J.~D. (1996). Are thought experiments just what you thought? \emph{Canadian Journal of Philosophy}, 26(3), 333--366.
\end{thebibliography}

\end{document}
